
\section{Advanced Activities}

The following activities are intended for strong master's or PhD students.
They emphasize constrained optimality, nonconvexity, numerical conditioning,
total derivatives, simulation-based optimization, and multidisciplinary
design optimization. MATLAB, Python, IPOPT, SNOPT, \texttt{fmincon},
OpenMDAO, or equivalent tools may be used where specified.

\subsection*{Problem 1: KKT Analysis of a Minimum-Mass Cantilever Beam}

A rectangular cantilever beam has width $b$, height $h$, length $L$,
elastic modulus $E$, material density $\rho$, and an end load $P$.
Its mass is
\begin{equation}
m(b,h)=\rho Lbh.
\end{equation}

The maximum bending stress and tip displacement are
\begin{equation}
\sigma(b,h)=\frac{6PL}{bh^2},
\end{equation}
and
\begin{equation}
\delta(b,h)=\frac{4PL^3}{Ebh^3}.
\end{equation}

The optimization problem is
\begin{align}
\min_{b,h}\quad
& \rho Lbh,
\\
\text{subject to}\quad
& \frac{6PL}{bh^2}-\sigma_{\max}\leq0,
\\
& \frac{4PL^3}{Ebh^3}-\delta_{\max}\leq0,
\\
& b_{\min}\leq b\leq b_{\max},
\\
& h_{\min}\leq h\leq h_{\max}.
\end{align}

Use
\begin{equation}
L=1.2~\mathrm{m},
\qquad
P=8~\mathrm{kN},
\qquad
E=70~\mathrm{GPa},
\end{equation}
\begin{equation}
\rho=2700~\mathrm{kg/m^3},
\qquad
\sigma_{\max}=120~\mathrm{MPa},
\qquad
\delta_{\max}=4~\mathrm{mm},
\end{equation}
and
\begin{equation}
0.01\leq b\leq0.12~\mathrm{m},
\qquad
0.02\leq h\leq0.20~\mathrm{m}.
\end{equation}

\begin{enumerate}
    \item Write the Lagrangian, including multipliers for the stress,
    displacement, and bound constraints.

    \item Write all KKT conditions: stationarity, primal feasibility,
    dual feasibility, and complementary slackness.

    \item Assume initially that the stress and displacement constraints are
    active and all bounds are inactive. Solve the two active constraints
    analytically for $b$ and $h$.

    \item Check whether the resulting design satisfies the variable bounds.

    \item Compute the associated Lagrange multipliers and determine whether
    the assumed active set satisfies dual feasibility.

    \item Enumerate all physically plausible active sets involving:
    \begin{enumerate}
        \item stress only,
        \item displacement only,
        \item stress and displacement,
        \item one performance constraint and one variable bound.
    \end{enumerate}

    \item Determine the globally optimal feasible design by comparing all
    valid KKT candidates.

    \item Interpret each nonzero multiplier as the local value of relaxing
    its corresponding engineering requirement.
\end{enumerate}


\subsection*{Problem 2: Local and Global Minima in a Tilted Double-Well Problem}

Consider the nonconvex function
\begin{equation}
f(x,y)
=
(x^2-1)^2
+
(y-x)^2
+
0.15x.
\end{equation}

\begin{enumerate}
    \item Derive the gradient and Hessian.

    \item Show that every stationary point satisfies
    \begin{equation}
    y=x
    \end{equation}
    and
    \begin{equation}
    4x^3-4x+0.15=0.
    \end{equation}

    \item Compute all real roots of the cubic equation to at least eight
    significant digits.

    \item Classify every stationary point using the Hessian eigenvalues.

    \item Identify the local minima, the saddle point, and the global minimum.

    \item Apply gradient descent with an Armijo backtracking line search from
    the initial points
    \begin{equation}
    (-2,2),\quad
    (-1,0),\quad
    (0,0),\quad
    (1,0),\quad
    (2,-2).
    \end{equation}

    \item Repeat using Newton's method with line search.

    \item Construct a basin-of-attraction plot on
    \begin{equation}
    -2\leq x\leq2,
    \qquad
    -2\leq y\leq2.
    \end{equation}

    \item Explain why convergence from several initial guesses provides
    evidence about nonconvexity but does not prove global optimality.

    \item Add the constraint
    \begin{equation}
    x+y\geq0
    \end{equation}
    and determine the constrained global minimum using KKT conditions.
\end{enumerate}


\subsection*{Problem 3: Scaling, Conditioning, and Gradient-Method Convergence}

Consider the quadratic optimization problem
\begin{equation}
\min_{\mathbf{x}\in\mathbb{R}^3}
\quad
f(\mathbf{x})
=
\frac{1}{2}\mathbf{x}^TH\mathbf{x}
-
\mathbf{b}^T\mathbf{x},
\end{equation}
where
\begin{equation}
H=
\begin{bmatrix}
10^8&0&0\\
0&1&0\\
0&0&10^{-4}
\end{bmatrix},
\qquad
\mathbf{b}=
\begin{bmatrix}
10^4\\
1\\
10^{-2}
\end{bmatrix}.
\end{equation}

\begin{enumerate}
    \item Compute the exact minimizer
    \begin{equation}
    \mathbf{x}^*=H^{-1}\mathbf{b}.
    \end{equation}

    \item Compute the $2$-norm condition number
    \begin{equation}
    \kappa_2(H).
    \end{equation}

    \item For fixed-step gradient descent,
    \begin{equation}
    \mathbf{x}^{(k+1)}
    =
    \mathbf{x}^{(k)}
    -
    \alpha
    \nabla f
    \left(
    \mathbf{x}^{(k)}
    \right),
    \end{equation}
    derive the largest stable step size.

    \item Derive the optimal fixed step size
    \begin{equation}
    \alpha^*
    =
    \frac{2}{\lambda_{\max}+\lambda_{\min}},
    \end{equation}
    and the corresponding worst-case error contraction factor
    \begin{equation}
    q^*
    =
    \frac{\kappa_2(H)-1}
    {\kappa_2(H)+1}.
    \end{equation}

    \item Estimate the number of iterations required to reduce the error norm
    by a factor of $10^{-6}$.

    \item Introduce the scaled variables
    \begin{equation}
    \mathbf{x}=S\mathbf{z},
    \qquad
    S=
    \begin{bmatrix}
    10^{-4}&0&0\\
    0&1&0\\
    0&0&10^2
    \end{bmatrix}.
    \end{equation}

    \item Derive the scaled Hessian
    \begin{equation}
    \widetilde{H}=S^THS
    \end{equation}
    and compute its condition number.

    \item Compare gradient descent, Newton's method, and BFGS on the original
    and scaled formulations from
    \begin{equation}
    \mathbf{x}^{(0)}=
    \begin{bmatrix}
    1\\1\\1
    \end{bmatrix}.
    \end{equation}

    \item Explain why variable scaling changes numerical behavior without
    changing the underlying physical optimum.

    \item Give one example of physical ill-conditioning that cannot be
    eliminated by simple coordinate scaling.
\end{enumerate}


\subsection*{Problem 4: Total Derivatives through a Coupled Multidisciplinary Model}

Consider two coupled disciplines:
\begin{align}
r_1
&=
y_1-x_1^2-\frac{1}{2}y_2
=
0,
\\
r_2
&=
y_2-\sin x_2-\frac{1}{4}y_1
=
0.
\end{align}

The objective and constraint are
\begin{equation}
f
=
(y_1-1)^2
+
(y_2-0.5)^2
+
0.1(x_1^2+x_2^2),
\end{equation}
and
\begin{equation}
g=y_1+y_2-2\leq0.
\end{equation}

The design bounds are
\begin{equation}
-1.5\leq x_1\leq1.5,
\qquad
-2\leq x_2\leq2.
\end{equation}

\begin{enumerate}
    \item Write the coupled residual system in the compact form
    \begin{equation}
    \mathbf{r}
    \left(
    \mathbf{x},
    \mathbf{y}
    \right)
    =
    \mathbf{0}.
    \end{equation}

    \item Solve the coupled equations analytically for
    \begin{equation}
    y_1(\mathbf{x}),
    \qquad
    y_2(\mathbf{x}).
    \end{equation}

    \item Derive the total state sensitivity using the implicit-function
    theorem:
    \begin{equation}
    \frac{d\mathbf{y}}{d\mathbf{x}}
    =
    -
    \left(
    \frac{\partial\mathbf{r}}{\partial\mathbf{y}}
    \right)^{-1}
    \frac{\partial\mathbf{r}}{\partial\mathbf{x}}.
    \end{equation}

    \item Derive the total derivatives
    \begin{equation}
    \frac{df}{d\mathbf{x}},
    \qquad
    \frac{dg}{d\mathbf{x}}.
    \end{equation}

    \item Evaluate the analytical total derivatives at
    \begin{equation}
    \mathbf{x}
    =
    \begin{bmatrix}
    0.8\\0.4
    \end{bmatrix}.
    \end{equation}

    \item Verify the derivatives using complex-step differentiation.

    \item Show that using only the partial derivative
    \begin{equation}
    \frac{\partial f}{\partial\mathbf{x}}
    \end{equation}
    while ignoring
    \begin{equation}
    \frac{d\mathbf{y}}{d\mathbf{x}}
    \end{equation}
    gives an incorrect search direction.

    \item Formulate and solve the problem using a multidisciplinary feasible
    architecture in which the coupling equations are converged at every
    design iteration.

    \item Reformulate the same problem using an individual-discipline
    feasible architecture in which $y_1$ and $y_2$ are optimization
    variables and the coupling equalities are consistency constraints.

    \item Compare the two formulations in terms of variable count,
    constraint count, derivative structure, and convergence.
\end{enumerate}


\subsection*{Problem 5: Discrete Adjoint Sensitivities for Simulation-Based Optimization}

Consider the nonlinear time-marching model
\begin{equation}
x_{k+1}
=
x_k
+
h
\left(
-p_1x_k
-p_2x_k^3
+u_k
\right),
\qquad
k=0,\ldots,N-1,
\end{equation}
with
\begin{equation}
x_0=1,
\qquad
h=0.05,
\qquad
N=100.
\end{equation}

The design variables are
\begin{equation}
\mathbf{z}
=
\begin{bmatrix}
p_1&
p_2&
u_0&
u_1&
\cdots&
u_{N-1}
\end{bmatrix}^{T}.
\end{equation}

Minimize
\begin{equation}
J
=
\frac{1}{2}
\left(
x_N-x_{\mathrm{target}}
\right)^2
+
\frac{h}{2}
\sum_{k=0}^{N-1}
\left(
q x_k^2
+
r u_k^2
\right)
+
\frac{\gamma}{2}
\left(
p_1^2+p_2^2
\right),
\end{equation}
where
\begin{equation}
x_{\mathrm{target}}=0,
\qquad
q=1,
\qquad
r=0.02,
\qquad
\gamma=0.01.
\end{equation}

\begin{enumerate}
    \item Derive the direct sensitivity recursion for
    \begin{equation}
    \frac{\partial x_k}{\partial z_j}.
    \end{equation}

    \item Determine the computational complexity of forming the complete
    gradient using direct sensitivities.

    \item Introduce a discrete adjoint sequence $\lambda_k$ and derive its
    backward recursion.

    \item Derive
    \begin{equation}
    \frac{\partial J}{\partial p_1},
    \qquad
    \frac{\partial J}{\partial p_2},
    \qquad
    \frac{\partial J}{\partial u_k}.
    \end{equation}

    \item Implement direct sensitivity, adjoint sensitivity, forward finite
    difference, central finite difference, and complex-step gradients.

    \item Perform a directional-derivative test using a random normalized
    direction $\mathbf{v}$:
    \begin{equation}
    \frac{
    J(\mathbf{z}+\epsilon\mathbf{v})
    -
    J(\mathbf{z})
    }{\epsilon}
    \approx
    \nabla J(\mathbf{z})^T\mathbf{v}.
    \end{equation}

    \item Compare gradient errors for
    \begin{equation}
    10^{-2}\leq\epsilon\leq10^{-14}.
    \end{equation}

    \item Compare computational cost as $N$ increases from
    \begin{equation}
    100
    \quad\text{to}\quad
    10\,000.
    \end{equation}

    \item Explain why the adjoint method is especially attractive when the
    number of design variables is large and the number of scalar outputs is
    small.

    \item Use an optimization solver to minimize $J$ subject to
    \begin{equation}
    0.1\leq p_1\leq3,
    \qquad
    0\leq p_2\leq2,
    \qquad
    -2\leq u_k\leq2.
    \end{equation}
\end{enumerate}


\subsection*{Problem 6: OpenMDAO Benchmark using the Sellar Coupled Problem}

Consider the classical two-discipline Sellar problem:
\begin{align}
y_1
&=
z_1^2+z_2+x-0.2y_2,
\\
y_2
&=
\sqrt{y_1}+z_1+z_2.
\end{align}

The optimization problem is
\begin{align}
\min_{x,z_1,z_2}\quad
&
f
=
x^2+z_2+y_1+e^{-y_2},
\\
\text{subject to}\quad
&
3.16-y_1\leq0,
\\
&
y_2-24\leq0,
\end{align}
with bounds
\begin{equation}
0\leq x\leq10,
\qquad
-10\leq z_1\leq10,
\qquad
0\leq z_2\leq10.
\end{equation}

\begin{enumerate}
    \item Identify:
    \begin{enumerate}
        \item local design variables,
        \item shared design variables,
        \item discipline outputs,
        \item coupling variables,
        \item system-level objective and constraints.
    \end{enumerate}

    \item Derive the coupled residual equations.

    \item Derive the fixed-point iteration that results from solving the
    disciplines in Gauss--Seidel order.

    \item Determine the local convergence condition for the coupling
    iteration from the spectral radius of the iteration Jacobian.

    \item Implement the multidisciplinary feasible formulation in OpenMDAO,
    using a nonlinear block Gauss--Seidel or Newton solver for the coupled
    analysis.

    \item Supply exact partial derivatives for both disciplines and use
    OpenMDAO's total-derivative check.

    \item Solve the optimization using at least three different initial
    designs.

    \item Reformulate the problem using an individual-discipline feasible
    architecture by promoting $y_1$ and $y_2$ to design variables and adding
    consistency constraints.

    \item Compare the two architectures in terms of:
    \begin{enumerate}
        \item number of design variables,
        \item number of constraints,
        \item coupling-solver iterations,
        \item optimizer iterations,
        \item total derivative accuracy,
        \item final objective value.
    \end{enumerate}

    \item Perturb the discipline equations by replacing
    \begin{equation}
    -0.2y_2
    \end{equation}
    with
    \begin{equation}
    -\beta y_2,
    \end{equation}
    where
    \begin{equation}
    0.1\leq\beta\leq0.8.
    \end{equation}
    Determine how coupling strength affects fixed-point convergence and
    optimization performance.

    \item Explain why a converged multidisciplinary analysis does not by
    itself prove that the optimization result is globally optimal or
    numerically reliable.
\end{enumerate}
